% `tex.web` §1105 `delete_last`: each of the three removers takes only its own
% kind of node.
%
%     if not is_char_node(tail) then
%       if type(tail)=cur_chr then …remove it…
%
% `cur_chr` is the node type the command names — glue for `\unskip`, a kern for
% `\unkern`, a penalty for `\unpenalty` — and if the last node is anything else
% nothing happens at all.
%
% We had one primitive for all three and removed whichever discardable happened
% to be last, so `\unskip` ate a kern and `\unpenalty` ate glue. That is how a
% document loses an italic correction it never asked to lose: `\/\unskip` is
% ordinary, and the kern `\/` just added is not what `\unskip` is for.
\catcode`\{=1 \catcode`\}=2
\showboxbreadth=100 \showboxdepth=100
\font\rip=cmr10 \rip

% Each remover against a node of the wrong kind: nothing goes.
\setbox0\hbox{A\kern3pt\unskip}
\showbox0
\setbox1\hbox{A\hskip3pt\unkern}
\showbox1
\setbox2\hbox{A\hskip3pt\unpenalty}
\showbox2

% ...and against its own kind: it goes.
\setbox3\hbox{A\hskip3pt\unskip}
\showbox3
\setbox4\hbox{A\kern3pt\unkern}
\showbox4
\setbox5\hbox{A\penalty7\unpenalty}
\showbox5
\end
